The matrix-unit relations are the hyperbolic-pair identities \eqref{eq:hyperbolicpair} read off directly: $E_{12}E_{21} = ff^* = E_{11}$, $E_{21}E_{12} = f^*f = E_{22}$, $E_{12}^2 = f^2 = 0 = E_{21}^2$, $E_{11}E_{12} = ff^*f = f = E_{12}$, $E_{12}E_{11} = f\,ff^* = 0$, and so on; $E_{11} + E_{22} = 1$ is the pair's anticommutator. For the corner: using the displayed decomposition of $B$ and $ff^* a = a\, ff^*$ (evenness), one computes $ff^* (a_0 + a_1 f + a_2 f^* + a_3 ff^*)\, ff^* = (a_0 + a_3) ff^*$, since $ff^* \cdot f \cdot ff^* = f f^* Q(f) f^* = f \cdot ff^* = f^2 f^* = 0$ and likewise the $f^*$ term vanishes, while the $1$ and $ff^*$ terms remain; so $E_{11}BE_{11} = A\,ff^*$. The map $a \mapsto a\,ff^*$ is an algebra isomorphism onto it: multiplicative since $(a\,ff^*)(b\,ff^*) = ab\,(ff^*)^2 = ab\,ff^*$, unital onto the corner unit $ff^*$, injective by the directness of the decomposition. The Peirce isomorphism of Lemma \claimref{transport.matrix-units-peirce-isomorphism}, whose proof uses only the matrix-unit relations, then gives $B \cong \Mat_2(A)$. The embedding: for $a \in A$, $\varphi(a)_{11} = E_{11} a E_{11} = a\,ff^*$ and $\varphi(a)_{22} = E_{12}\, a\, E_{21} = f a f^* = \hat{a}\, ff^*$, while the off-diagonal entries vanish ($E_{11} a E_{21} = a\, ff^*f^* = 0$, $E_{12} a E_{11} = \hat a\,f ff^* = 0$); under the corner identification this is $\mathrm{diag}(a, \hat{a})$, which is multiplicative since the grade involution is an automorphism (\claimref{std-algebra.involution-properties}). Morita-class preservation, and with it the invariance of the type-C diagnosis, is the statement $\Mat_2(A) \sim_{\text{Morita}} A$ (\claimref{std-algebra.morita-amplification}) together with the fact that the Wedderburn data of $\Mat_2(\Mat_m(D))$ (or of the split case, blockwise) is that of $A$ with $m$ doubled.
